% =============================================================================
%  ch08_boundary_conditions.tex — Generic Boundary \& Initial Condition Design
% =============================================================================
\chapter{Generic Boundary and Initial Condition Infrastructure}
\label{ch:boundary-conditions}

This chapter documents the design and implementation of the generic
boundary condition (BC) and initial condition (IC) infrastructure for
\texttt{fall\_n}.  The goals are:
\begin{enumerate}[nosep]
  \item \textbf{Per-DOF granularity:}  constrain individual degrees of freedom
        (roller supports, symmetry BCs) rather than all DOFs simultaneously.
  \item \textbf{Imposed (non-zero) Dirichlet:}  prescribe non-zero
        displacement, temperature, or any nodal field value.
  \item \textbf{Compile-time node selectors:}  apply constraints via
        geometric predicates (plane, box, sphere, custom lambda) that
        resolve at compile time with zero runtime overhead.
  \item \textbf{State transfer:}  capture the complete post-solve state
        (\texttt{ModelState}) and use it as the initial condition for
        a subsequent analysis---enabling staged construction, restart,
        and multi-physics chaining.
  \item \textbf{Extensibility:}  the design should accommodate future
        field types (temperature, initial stress, pore pressure) without
        modifying the selector or constraint machinery.
\end{enumerate}


% ═════════════════════════════════════════════════════════════════════════
\section{Gap Analysis: the Legacy API}
\label{sec:bc-gap}

The original Model API exposed these methods for boundary conditions:

\begin{lstlisting}
void fix_node(size_t node_idx);                          // ALL DOFs = 0
void fix_x(double val, double tol = 1e-6);               // all DOFs at x=val
void apply_node_force(size_t node_idx, auto... f);        // nodal Neumann
void apply_surface_traction(string group, auto... t);     // consistent traction
\end{lstlisting}

\begin{table}[ht]
\centering\small
\caption{Limitations of the legacy BC interface.}
\label{tab:bc-gaps}
\begin{tabular}{lp{8cm}}
\toprule
\textbf{Limitation} & \textbf{Impact} \\
\midrule
\texttt{fix\_node} constrains \emph{all} DOFs &
Cannot model roller supports, symmetry planes, or mixed BCs. \\
\texttt{fix\_x/y/z} misnamed &
These fix all three displacement components, not just the named one.
Misleading to users.\\
No non-zero Dirichlet via Model API &
\texttt{global\_imposed\_solution\_} was created and zeroed but never
populated through the Model interface.  Only
\texttt{DynamicAnalysis::PrescribedBC} filled it.\\
\texttt{ElementInitialStress} unused &
The struct existed in \texttt{BoundaryCondition.hh} and was storable in
\texttt{BoundaryConditionSet}, but no evaluator integrated it into the
assembly pipeline.\\
Procedural selection &
Every \texttt{fix\_*} method hard-coded a coordinate loop.  No reusable
geometric predicate.\\
\texttt{VecAssemblyBegin/End} per call &
\texttt{apply\_node\_force} triggered PETSc assembly for each node. \\
\bottomrule
\end{tabular}
\end{table}


% ═════════════════════════════════════════════════════════════════════════
\section{Design Decisions and Trade-offs}
\label{sec:bc-design}

\subsection{Compile-Time Polymorphism via Concepts}

Consistent with \texttt{fall\_n}'s policy-based architecture
(\Cref{ch:architecture}), the node selector system uses C++20 concepts
rather than virtual dispatch:

\begin{lstlisting}
namespace fn {
template <typename S, typename NodeT>
concept NodeSelector = requires(const S& s, const NodeT& n) {
    { s(n) } -> std::convertible_to<bool>;
};
}
\end{lstlisting}

\noindent
Any callable \lstinline{f(const Node<dim>&) -> bool} satisfies
\texttt{NodeSelector}: lambdas, function objects, and the built-in
selectors all qualify.  The compiler resolves the predicate at compile
time, and after inlining a \texttt{PlaneSelector} collapses to a single
\texttt{fabs} + comparison---zero overhead compared to a hand-written
loop.

\paragraph{Why not \texttt{std::function}?}
A \texttt{std::function<bool(const Node\&)>} would erase the type and
force heap allocation for lambdas that exceed the small-buffer size (typically
16--32 bytes).  For selectors applied to every node in the mesh, this
per-call overhead is measurable.  The concept-based approach pays nothing
at runtime.

\paragraph{Why not virtual?}
A virtual \texttt{AbstractSelector::operator()} would add one indirect
call per node per constraint pass.  With 10\,000 nodes this is
$\sim$10\,$\mu$s (branch-predictor penalty), acceptable but unnecessary
when the concrete type is known at compile time.


\subsection{Constraint Storage: Extending the Existing Map}

The existing \texttt{ConstraintDofInfo} type already maps sieve point IDs
to (DOF index list, prescribed value list):

\begin{lstlisting}
using ConstraintDofInfo =
  std::map<PetscInt,
           std::pair<std::vector<PetscInt>,
                     std::vector<PetscScalar>>>;
\end{lstlisting}

\noindent
The new API \emph{extends} this map rather than replacing it.  The key
operations are:

\begin{enumerate}[nosep]
  \item \textbf{Insert in sorted order.}  PETSc's
        \texttt{PetscSectionSetConstraintIndices} requires increasing
        DOF indices.  The \texttt{constrain\_dof} method uses
        \texttt{std::lower\_bound} to maintain this invariant.
  \item \textbf{Merge, not overwrite.}  If DOF~0 of a node is already
        constrained and the user constrains DOF~2, both are stored.
        An update to DOF~0 changes the prescribed value without
        duplicating the index.
\end{enumerate}

\subsection{Non-Zero Dirichlet: Where Values Flow}
\label{sec:bc-nonzero}

Non-zero prescribed values are stored in \texttt{constraints\_} alongside
the DOF indices.  During \texttt{setup()}, a new
\texttt{fill\_imposed\_solution()} step extracts these values from the
constraint map into the PETSc local vector
\texttt{global\_imposed\_solution\_}:

\begin{lstlisting}
void setup() {
    setup_boundary_conditions();  // PETSc section
    setup_vectors();              // create Vecs
    fill_imposed_solution();      // NEW: u_D -> imposed
    setup_matrix();               // create Mat
}
\end{lstlisting}

\begin{remark}
Non-zero Dirichlet requires that the analysis accounts for the coupling
$\mathbf{K}_{fc}\,\mathbf{u}_c$ in the right-hand side.
\texttt{NonlinearAnalysis} handles this naturally: the internal forces
$\mathbf{f}_{\text{int}}(\mathbf{u}_{\text{free}} + \mathbf{u}_D)$ include
the prescribed part.  \texttt{LinearAnalysis} currently does \emph{not}
apply this correction---for non-zero Dirichlet, users should solve with
\texttt{NonlinearAnalysis} (which converges in one Newton step for linear
problems).
\end{remark}


\subsection{ModelState: Platform-Independent Snapshot}

For state transfer between models, we need a data structure that:
\begin{itemize}[nosep]
  \item Does \emph{not} depend on PETSc (no \texttt{Vec}, \texttt{Mat}).
  \item Does \emph{not} depend on the \texttt{MaterialPolicy} template parameter
        (so an elastic model's state can seed a plastic model).
  \item Is copyable, movable, and (future) serialisable.
\end{itemize}

\texttt{ModelState<dim>} satisfies all three.  It stores:
\begin{itemize}[nosep]
  \item Per-node: ID, displacement, velocity.
  \item Per-element, per-Gauss-point: stress and strain in Voigt notation
        (\texttt{std::array<double, nvoigt>}).
\end{itemize}

The PETSc-to-array extraction in \texttt{capture\_state()} uses
\texttt{VecGetArrayRead} (zero-copy read) and the element's
\texttt{material\_points()} accessor which recomputes stress from the
committed strain via the constitutive relation.


% ═════════════════════════════════════════════════════════════════════════
\section{Implementation}
\label{sec:bc-impl}

\subsection{NodeSelector.hh — Built-in Selectors}

\begin{table}[ht]
\centering\small
\caption{Built-in node selectors.}
\label{tab:selectors}
\begin{tabular}{lll}
\toprule
\textbf{Selector} & \textbf{Predicate} & \textbf{Parameters} \\
\midrule
\texttt{PlaneSelector<dim>}  & $|x_i - v| < \varepsilon$
                             & axis $i$, value $v$, tolerance $\varepsilon$ \\
\texttt{BoxSelector<dim>}    & $x \in [\text{lo},\text{hi}]$
                             & corners, tolerance \\
\texttt{SphereSelector<dim>} & $\|x - c\| \le r$
                             & centre $c$, radius $r$, tolerance \\
\texttt{NodeIdSelector}      & $\text{id} \in S$
                             & set of node IDs \\
\bottomrule
\end{tabular}
\end{table}

\noindent
Combinators turn simple selectors into compound predicates:

\begin{lstlisting}
auto corner = fn::select_and(
    fn::PlaneSelector<3>{0, 0.0},   // x = 0
    fn::PlaneSelector<3>{1, 0.0}    // y = 0
);
auto not_x0 = fn::select_not(fn::PlaneSelector<3>{0, 0.0});
auto union_ = fn::select_or(selectorA, selectorB);
\end{lstlisting}

\noindent
Convenience factories (\texttt{fn::on\_plane}, \texttt{fn::in\_box},
\texttt{fn::in\_sphere}, \texttt{fn::node\_ids}) provide a terser syntax.


\subsection{Model.hh — New Constraint Methods}

\begin{table}[ht]
\centering\small
\caption{New Model constraint API.}
\label{tab:constraint-api}
\begin{tabular}{lp{7.5cm}}
\toprule
\textbf{Method} & \textbf{Description} \\
\midrule
\texttt{constrain\_dof(n, d, v)}
  & Constrain DOF~$d$ of node~$n$ to value~$v$.  Sorted insertion. \\
\texttt{constrain\_node(n, vals)}
  & Constrain all DOFs of node~$n$ to the given array of values. \\
\texttt{constrain\_dof\_where(sel, d, v)}
  & Constrain DOF~$d$ at every node matching selector \texttt{sel}. \\
\texttt{constrain\_where(sel)}
  & Constrain all DOFs (homogeneous) at every matching node. \\
\texttt{constrain\_where(sel, vals)}
  & Constrain all DOFs with prescribed values at every matching node. \\
\midrule
\texttt{is\_constrained(n)}
  & Query: is any DOF of node~$n$ constrained? \\
\texttt{is\_dof\_constrained(n, d)}
  & Query: is DOF~$d$ of node~$n$ constrained? \\
\texttt{prescribed\_value(n, d)}
  & Get the prescribed displacement at DOF~$d$ (0 if free). \\
\texttt{num\_constraints()}
  & Total number of constrained DOFs across all nodes. \\
\midrule
\texttt{capture\_state()}
  & Extract current solution into a \texttt{ModelState<dim>}. \\
\texttt{apply\_initial\_displacement(s)}
  & Populate imposed solution from a \texttt{ModelState<dim>}. \\
\bottomrule
\end{tabular}
\end{table}


\subsection{Backward Compatibility}

The legacy methods \texttt{fix\_node}, \texttt{fix\_x}, \texttt{fix\_y},
\texttt{fix\_z}, and \texttt{fix\_orthogonal\_plane} remain unchanged in
semantics---they continue to constrain \emph{all} DOFs of every matching
node.  Internally, \texttt{fix\_orthogonal\_plane} now delegates to:
\begin{lstlisting}
void fix_orthogonal_plane(int i, double val, double tol) {
    constrain_where(fn::PlaneSelector<dim>{i, val, tol});
}
\end{lstlisting}
The code path is identical; only the implementation reuses the selector
infrastructure.


% ═════════════════════════════════════════════════════════════════════════
\section{Use-Case Walkthrough: Roller vs.\ Clamped Cantilever}
\label{sec:bc-roller}

Consider a cantilever beam $[0,L]\times[0,b]\times[0,h]$ loaded at
$x = L$ by a surface traction $\mathbf{t} = (0, t_y, t_z)$.

\paragraph{Clamped support (legacy):}
\begin{lstlisting}
M.fix_x(0.0);   // constrains u_x, u_y, u_z at x=0
\end{lstlisting}
Equivalent new API:
\begin{lstlisting}
M.constrain_where(fn::PlaneSelector<3>{0, 0.0});
\end{lstlisting}

\paragraph{Roller support (new API):}
Only $u_x = 0$ at $x = 0$.  Rigid body prevention requires additionally
constraining $u_y$ and $u_z$ at isolated points:
\begin{lstlisting}
// u_x = 0 on entire face
M.constrain_dof_where(fn::PlaneSelector<3>{0, 0.0}, 0, 0.0);

// Prevent rigid translation + rotation about x:
auto origin = fn::select_and(fn::PlaneSelector<3>{0, 0.0},
    fn::select_and(fn::PlaneSelector<3>{1, 0.0},
                   fn::PlaneSelector<3>{2, 0.0}));
M.constrain_dof_where(origin, 1, 0.0);  // u_y at origin
M.constrain_dof_where(origin, 2, 0.0);  // u_z at origin

auto top_z = fn::select_and(fn::PlaneSelector<3>{0, 0.0},
    fn::select_and(fn::PlaneSelector<3>{1, 0.0},
                   fn::PlaneSelector<3>{2, h}));
M.constrain_dof_where(top_z, 1, 0.0);   // prevents Rx
\end{lstlisting}

The roller case has 30~constrained DOFs (28~nodes $\times$ 1 + 2~extra)
versus 84 (28 $\times$ 3) for the clamped case.  With fewer constraints,
the beam is more flexible:

\begin{table}[ht]
\centering\small
\caption{Roller vs.\ clamped comparison (HEX8 beam, $E=200$, $\nu=0.3$).}
\label{tab:roller-clamped}
\begin{tabular}{lccc}
\toprule
\textbf{BC type} & \textbf{Constrained DOFs} & $\max|u|$ & \textbf{Ratio} \\
\midrule
Clamped & 84 & 0.594 & 1.00 \\
Roller ($u_x$ only) & 30 & 2.762 & 4.64$\times$ \\
\bottomrule
\end{tabular}
\end{table}

\noindent
This confirms: the per-DOF mechanism works correctly and produces a
physically meaningful stiffer/more-flexible response.


% ═════════════════════════════════════════════════════════════════════════
\section{State Capture and Continuation Analysis}
\label{sec:bc-state}

\subsection{ModelState Structure}

\begin{lstlisting}
template <size_t dim>
struct ModelState {
    static constexpr size_t nvoigt = dim * (dim + 1) / 2;

    struct NodalData {
        size_t node_id;
        array<double, dim> displacement;
        array<double, dim> velocity;       // zero for static
    };

    struct GaussPointData {
        array<double, nvoigt> stress;
        array<double, nvoigt> strain;
    };

    struct ElementData {
        size_t element_index;
        vector<GaussPointData> gauss_points;
    };

    vector<NodalData>   nodes;
    vector<ElementData> elements;
};
\end{lstlisting}

\subsection{Capture Workflow}

\texttt{Model::capture\_state()} extracts:
\begin{enumerate}[nosep]
  \item \textbf{Nodal displacements:}  from the PETSc local vector
        \texttt{current\_state\_} via \texttt{VecGetArrayRead} (zero-copy).
        Each node's DOF offsets from \texttt{node.dof\_index()} are used to
        index into the raw array.
  \item \textbf{Element stress/strain:}  from each \texttt{MaterialPoint}'s
        committed strain (\texttt{current\_state()}) and recomputed stress
        (\texttt{compute\_response()}).  The Voigt vectors are copied
        component-by-component into \texttt{std::array<double, nvoigt>},
        guarded by the compile-time \texttt{num\_components} constant.
\end{enumerate}

\subsection{Future: Initial Stress Integration}

The existing \texttt{ElementInitialStress<dim>} struct in
\texttt{BoundaryCondition.hh} can be populated from
\texttt{ModelState::ElementData}:

\begin{lstlisting}
for (const auto& ed : state.elements) {
    ElementInitialStress<dim> eis;
    eis.element_index = ed.element_index;
    for (const auto& gp : ed.gauss_points) {
        eis.gauss_point_stresses.push_back(gp.stress);
    }
    bcs.add_initial_stress(std::move(eis));
}
\end{lstlisting}

\noindent
For this to take effect, the element assembly must incorporate
$\bsigma_0$ into the internal force:
\[
  \mathbf{f}_{\text{int}}(\mathbf{u})
  = \int_{\Omega_e} \mathbf{B}^T \bigl(\bsigma(\bvarepsilon(\mathbf{u}))
      + \bsigma_0\bigr)\,\dd V
\]
This requires modifications to \texttt{ContinuumElement::compute\_internal\-\_force\_vector},
specifically adding a \texttt{$\sigma$\_initial} term at each Gauss point.
This integration is deferred to a future phase.


% ═════════════════════════════════════════════════════════════════════════
\section{Comparison with OpenSees and Kratos}
\label{sec:bc-comparison}

\begin{table}[ht]
\centering\small
\caption{BC/IC handling comparison across frameworks.}
\label{tab:bc-framework}
\begin{tabular}{p{3cm}p{3.5cm}p{3.5cm}p{3.5cm}}
\toprule
\textbf{Feature} & \textbf{OpenSees} & \textbf{Kratos} & \textbf{fall\_n} \\
\midrule
Single-DOF fix
  & \texttt{fix nodeTag dof}
  & \texttt{DISPLACEMENT\_X: 0.0}
  & \texttt{constrain\_dof(n,d,v)} \\
Geometric selection
  & Tcl loops
  & Python \texttt{FindNodalH}
  & \texttt{PlaneSelector}, lambda \\
Non-zero Dirichlet
  & \texttt{imposedMotion} pattern
  & \texttt{ApplyDirichlet\-Process}
  & \texttt{constrain\_dof\_where(s,d,v)} \\
State transfer
  & \texttt{-database} serialise
  & \texttt{ModelPartIO}
  & \texttt{ModelState} capture/apply \\
Initial stress
  & Per-element update
  & \texttt{INITIAL\_STRESS\_TENSOR}
  & Defined, not yet integrated \\
Dispatch cost
  & Virtual (runtime)
  & Virtual + Python (runtime)
  & Compile-time (zero) \\
\bottomrule
\end{tabular}
\end{table}

\paragraph{Key differences:}
\begin{itemize}[nosep]
  \item \textbf{OpenSees} uses a tag-based Tcl scripting layer for BC
        specification.  Geometric loops require user scripting.  State transfer
        uses a persistence layer (\texttt{FE\_Datastore}).
  \item \textbf{Kratos} uses Python processes (\texttt{ApplyDirichletProcess},
        \texttt{ApplyNeumannProcess}) that iterate over model parts.
        Highly flexible but with Python $\to$ C++ boundary overhead.
  \item \textbf{fall\_n} resolves all selector logic at compile time.
        The constraint map lives in pure C++ with PETSc-section-based
        enforcement.  No scripting layer, no virtual dispatch per node.
\end{itemize}


% ═════════════════════════════════════════════════════════════════════════
\section{Extensibility Roadmap}
\label{sec:bc-extensibility}

The current infrastructure is designed for forward extension.
\Cref{tab:bc-roadmap} outlines planned capabilities and their
implementation path.

\begin{table}[ht]
\centering\small
\caption{Extensibility roadmap.}
\label{tab:bc-roadmap}
\begin{tabular}{p{3.5cm}p{5cm}p{5cm}}
\toprule
\textbf{Capability} & \textbf{Implementation path} & \textbf{Status} \\
\midrule
Per-DOF constraints
  & \texttt{constrain\_dof}, \texttt{constrain\_dof\_where}
  & \textbf{Implemented} \\
Non-zero Dirichlet
  & \texttt{fill\_imposed\_solution()} + NL correction
  & \textbf{Implemented} (NL only) \\
Node selectors
  & \texttt{NodeSelector} concept + built-in selectors
  & \textbf{Implemented} \\
State capture/transfer
  & \texttt{ModelState} + \texttt{capture\_state}
  & \textbf{Implemented} \\
\addlinespace
Initial stress
  & Extend \texttt{ContinuumElement::compute\_internal\_force}
    to add $\bsigma_0$ at each GP
  & Designed, deferred \\
Body forces
  & Consistent nodal load via $\int N^T\,\mathbf{b}\,\dd V$
  & Planned \\
Temperature field
  & Additional \texttt{ModelState} field + thermal strain
    $\bvarepsilon_\theta = \alpha\,\Delta T\,\mathbf{1}$
  & Planned \\
Multi-step builder
  & Fluent builder with chained \texttt{.constrain().material().build()}
  & Planned \\
State serialisation
  & HDF5/JSON export of \texttt{ModelState}
  & Planned \\
Linear non-zero BC
  & RHS correction $\mathbf{f}_{\text{eff}} = \mathbf{f}_{\text{ext}}
    - \mathbf{K}_{fc}\,\mathbf{u}_D$  in \texttt{LinearAnalysis}
  & Planned \\
\bottomrule
\end{tabular}
\end{table}


% ═════════════════════════════════════════════════════════════════════════
\section{Test Suite}
\label{sec:bc-tests}

The implementation is validated by 81 tests in
\texttt{tests/test\_boundary\_conditions.cpp}, organized as follows:

\begin{table}[ht]
\centering\small
\caption{Test suite structure (81 tests).}
\label{tab:bc-tests}
\begin{tabular}{lrl}
\toprule
\textbf{Category} & \textbf{Count} & \textbf{Coverage} \\
\midrule
Node selectors (pure geometry) & 38 &
  PlaneSelector, BoxSelector, SphereSelector, NodeIdSelector,
  combinators, lambdas, 2D/3D, factories \\
ModelState data structure & 9 &
  Default state, add/query, copy semantics \\
Per-DOF constraints & 11 &
  Single DOF, multiple DOFs, non-zero value, value update \\
Selector-based constraints & 10 &
  PlaneSelector on DOF, all DOFs, values, lambda, composite \\
Full solve: roller vs.\ clamped & 3 &
  Clamped $>$ 0, Roller $>$ 0, Roller $>$ Clamped \\
State capture after solve & 7 &
  Non-empty, node count, element count, non-zero
  displacement, non-zero stress, ID match \\
Backward compatibility & 3 &
  \texttt{fix\_x} gives positive deflection, reasonable magnitude \\
\bottomrule
\end{tabular}
\end{table}


% ═════════════════════════════════════════════════════════════════════════
\section{Summary}
\label{sec:bc-summary}

The generic BC/IC infrastructure extends \texttt{fall\_n} from an
all-or-nothing clamped-support tool to a framework capable of:

\begin{itemize}[nosep]
  \item \textbf{Arbitrary Dirichlet BCs:}  per-DOF, per-node, per-selector,
        with prescribed values (zero or non-zero).
  \item \textbf{Zero-cost selectors:}  compile-time geometric predicates
        with no runtime overhead versus hand-written loops.
  \item \textbf{State snapshots:}  PETSc-independent \texttt{ModelState}
        for continuation, restart, and future serialisation.
  \item \textbf{Full backward compatibility:}  \texttt{fix\_x}, \texttt{fix\_y},
        \texttt{fix\_z} continue to work exactly as before.
\end{itemize}

\noindent
The design follows the same policy-based, compile-time polymorphic
philosophy established in the element architecture (\Cref{ch:architecture})
and validated by an 81-test suite.
